\chapter{Parcours opérateur}

\section{Responsabilité}

L'opérateur supervise les ressources de son organisation. Il peut gérer les
stations et connecteurs, exécuter des commandes OCPP autorisées, traiter les
alertes, planifier la maintenance et échanger des rapports. Il n'administre pas
les organisations ni les comptes globaux de la plateforme.

\manualscreenshot{operator-overview.png}{Tableau de bord opérationnel}{fig:user-operator-overview}

\section{Superviser les stations}

La page \textbf{Stations} regroupe l'état calculé, la connectivité, la capacité
des connecteurs, les erreurs et la télémétrie. L'opérateur peut utiliser une vue
carte pour la localisation et une vue tableau pour comparer les états.

La fiche station sépare :

\begin{itemize}[leftmargin=1.6em]
  \item le résumé opérationnel ;
  \item la télémétrie récente ;
  \item les connecteurs et leur état OCPP ;
  \item l'historique des commandes ;
  \item les maintenances, sessions, alertes et documents.
\end{itemize}

\subsection{Mettre en service une station simulée}

L'action \textbf{Commission station} crée ensemble la station, ses connecteurs
physiques et son identité OCPP. Pour la cible \textbf{Local SAP simulator},
l'enregistrement du profil est automatisé par la plateforme : l'opérateur ne
copie pas de commande dans un terminal.

\begin{procedurebox}{Créer une station de simulation}
  \begin{enumerate}[leftmargin=1.6em]
    \item ouvrir \textbf{Stations}, puis \textbf{Commission station} ;
    \item saisir l'identité, l'adresse et choisir la position sur la carte ;
    \item définir le matériel et les connecteurs OCPP ;
    \item sélectionner la cible de simulation ;
    \item vérifier le récapitulatif et créer la station ;
    \item ouvrir sa fiche ou le \textbf{Simulation Lab} pour suivre connexion,
    démarrage et heartbeats.
  \end{enumerate}
\end{procedurebox}

\expected{La station apparaît avec un profil de simulateur traçable. Sa
disponibilité reste calculée à partir des événements reçus ; la création en base
ne suffit pas à la déclarer en ligne.}

\section{Commandes distantes}

\begin{longtable}{|L{3.4cm}|L{5.1cm}|L{7cm}|}
  \hline
  \rowcolor{ctmint}\textbf{Commande} & \textbf{But} & \textbf{Précaution} \\ \hline
  \endfirsthead
  \hline
  \rowcolor{ctmint}\textbf{Commande} & \textbf{But} & \textbf{Précaution} \\ \hline
  \endhead
  Restart station & Redémarrer le logiciel de la borne lorsqu'elle répond encore
  au canal OCPP. & Vérifier les sessions actives et privilégier un redémarrage
  doux avant un redémarrage forcé. \\ \hline
  Unlock connector & Demander le déverrouillage d'un câble. & Confirmer l'absence
  de charge active et la sécurité de l'utilisateur sur place. \\ \hline
  Set maintenance mode & Retirer temporairement la station du service normal. &
  Créer ou relier une maintenance et informer les acteurs concernés. \\ \hline
  Remote start / stop & Démarrer ou arrêter une transaction autorisée. & Vérifier
  l'identité du client, le connecteur et la raison opérationnelle. \\ \hline
\end{longtable}

\expected{Chaque demande génère une trace avec son état : en attente, envoyée,
acceptée, rejetée, expirée ou échouée. Une commande acceptée ne garantit pas à
elle seule le résultat physique ; les événements suivants doivent le confirmer.}

\section{Piloter le laboratoire OCPP}

Depuis une station simulée, \textbf{Open Simulation Lab} ouvre un espace dédié.
L'opérateur choisit une station et observe le lien OCPP, les connecteurs, les
impulsions en direct et l'historique des actions.

\begin{itemize}[leftmargin=1.6em]
  \item \textbf{Connect / Disconnect} teste la présence ou la perte du lien ;
  \item \textbf{Heartbeat} vérifie l'actualisation du dernier signal ;
  \item \textbf{Plug / Unplug} reproduit le branchement du câble ;
  \item \textbf{Inject fault / Recover} teste le défaut et son rétablissement ;
  \item \textbf{Cable cycle} et \textbf{Fault and recovery} exécutent des
  scénarios déterministes ;
  \item \textbf{Soft restart}, \textbf{Unlock} et \textbf{Maintenance mode}
  restent des commandes centrales, séparées de la simulation physique.
\end{itemize}

\attention{Déconnecter une station ou injecter un défaut peut créer une alerte
et modifier sa disponibilité. Utiliser ces actions sur l'environnement simulé
prévu, puis vérifier le retour à un état cohérent.}

\section{Traiter une alerte}

\manualscreenshot{operator-alerts.png}{Espace de traitement des alertes}{fig:user-operator-alerts}

\begin{procedurebox}{Prise en charge d'une alerte}
  \begin{enumerate}[leftmargin=1.6em]
    \item filtrer les alertes nouvelles ou critiques ;
    \item ouvrir l'alerte et vérifier station, connecteur, origine et horodatage ;
    \item reconnaître la prise en charge ;
    \item effectuer une action distante lorsque le diagnostic le permet ;
    \item créer une intervention si un déplacement est nécessaire ;
    \item assigner le technicien, la priorité et l'échéance ;
    \item suivre les preuves et résoudre l'alerte lorsque la cause est levée.
  \end{enumerate}
\end{procedurebox}

Une alerte décrit un événement à analyser. Une intervention est le travail
assigné pour le traiter. Une maintenance représente une période planifiée ou
coordonnée pendant laquelle un équipement peut être indisponible.

\section{Planifier la maintenance}

\manualscreenshot{operator-maintenance.png}{Planification des maintenances}{fig:user-operator-maintenance}

L'opérateur peut consulter le tableau ou le calendrier. Une maintenance précise
la station, le type, la période, le responsable et le statut. Les changements de
date doivent tenir compte des sessions prévues et des interventions déjà
assignées.

\begin{procedurebox}{Créer une maintenance}
  \begin{enumerate}[leftmargin=1.6em]
    \item sélectionner \textbf{Plan maintenance} ;
    \item choisir la station et, si nécessaire, le connecteur ;
    \item définir le type, la date de début et la durée ;
    \item assigner un responsable ;
    \item enregistrer puis vérifier le calendrier ;
    \item suivre le passage de planifiée à en cours puis terminée.
  \end{enumerate}
\end{procedurebox}

\section{Sessions, paiements et rapports}

L'opérateur consulte les sessions de son organisation et les paiements associés,
sans accéder aux données d'une autre organisation. L'espace \textbf{Reports}
sert à préparer une relève, documenter les incidents et transmettre un rapport
à un collègue ou à l'administrateur.

\manualscreenshot{operator-reports.png}{Rapports et relève de l'opérateur}{fig:user-operator-reports}

Un rapport peut être conservé en brouillon, envoyé, reçu ou archivé. La
composition propose des modèles adaptés au rôle, un destinataire, une catégorie,
une priorité, une période et jusqu'à cinq pièces jointes de 10 Mo chacune (PDF,
image, Word ou Excel). Si une pièce jointe échoue, le rapport reste enregistré
en brouillon afin de permettre une nouvelle tentative. Les pièces jointes ne
doivent jamais contenir de secrets, identifiants de paiement ou mots de passe
OCPP.
